% from ega4-18-fr.tex
\oldpage[IV-4]{202}

\begin{enumerate}
  \item[(i)] L'image de $X'$ par l'immersion fermée $j:X'\to\mathbf P(\mathscr E)$ correspondant à l'homomorphisme surjectif de $\sh{O}_X$-Algèbres graduées $\mathbf S_{\sh{O}_X}(\mathscr E)\to\mathbf S_{\sh{O}_X}(\mathscr J)\to\bigoplus_{n\geq0}\mathscr J^n$ (II, 3.6.2) est un sous-préschéma fermé de $Z$.
  \item[(ii)] Si $\mathscr E$ est un $\sh{O}_X$-Module localement libre de rang $m$, $\mathscr H$ est un $\sh{O}_P$-Module localement libre de rang $m-1$.
  \item[(iii)] On suppose que $X$ soit localement noethérien, que $\mathscr E$ soit localement libre de type fini, que le sous-préschéma $Y$ soit régulier, et qu'en tout point $x\in Y$, l'immersion canonique $i:Y\to X$ soit régulière et de codimension égale à $\operatorname{rg}_{\sh{O}_{X,x}}(\mathscr E_x)$ (ce que nous exprimerons en disant que l'homomorphisme $u$ est \emph{régulier} au point $x$ (cf. chap. V)). Alors l'image de $X'$ par l'immersion fermée $j:X'\to P$ est le sous-préschéma fermé $Z$ de $P$; en tous les points de $X'-Y'$ ($Y'$ image réciproque de $Y$ dans $X'$) l'immersion $j$ est transversalement régulière relativement à $X$; enfin en tous les points $x'$ de $Y'$, $X'$ est régulier, l'immersion $j$ est régulière et de codimension $\operatorname{rg}_{\sh{O}_{P,z}}(\mathscr H_z)$ où $z=j(x')$ (autrement dit l'homomorphisme $v:\mathscr H\to\sh{O}_P$ est régulier en $z$).
\end{enumerate}
\end{proposition}

Toutes les questions étant locales sur $X$, on peut se borner au cas où $X=\Spec(A)$ est affine et $\mathscr E=\widetilde E$, où $E$ est un $A$-module. Pour faire la vérification de (i), on peut en outre se borner à examiner ce qui se passe dans un ouvert affine de $P$ de la forme $D_+(t)$, où $t\in E=\mathbf S_A^1(E)$ (II, 2.3.14). Si l'on pose $S=\mathbf S_A(E)$, on a alors $\Gamma(D_+(t),\sh{O}_P)=S_{(t)}$ (II, 2.4.1), et $p^*(\mathscr E)|D_+(t)=(E\otimes_A S_{(t)})^\sim$. Si l'on se reporte à la définition de $\alpha_1^{\#}$ (II, 2.6.2.2), on voit qu'à tout élément
\[
  a=\sum_i x^{(i)}\otimes\frac{x_1^{(i)}x_2^{(i)}\cdots x_n^{(i)}}{t^n}
\]
de $E\otimes_A S_{(t)}$ (où $x^{(i)}$ et les $x_j^{(i)}$ sont dans $E$) $\alpha_1^{\#}$ fait correspondre l'élément
\[
  \sum_i\frac{x^{(i)}x_1^{(i)}\cdots x_n^{(i)}}{t^n}
\]
de $(S(1))_{(t)}$. Il s'agit (pour prouver (i)) de montrer que si ce dernier élément est nul, alors il en est de même de l'image de
\[
  a'=v(a)=\sum_i u(x^{(i)})\frac{x_1^{(i)}\cdots x_n^{(i)}}{t^n}
\]
par l'homomorphisme canonique $S_{(t)}\to(\bigoplus_{n\geq0}(u(E))^n)_{(u(t))}$ (II, 3.6.2). Or, cette image n'est autre que
\[
  \sum_i\frac{u(x^{(i)})u(x_1^{(i)})\cdots u(x_n^{(i)})}{(u(t))^n}
\]
dans l'anneau $A_{u(t)}$, c'est-à-dire le produit de l'élément $u(t)$ et de l'image canonique de $\sum_i x^{(i)}x_1^{(i)}\cdots x_n^{(i)}/t^{n+1}$ par l'homomorphisme d'algèbres, prolongement canonique à $S_{(t)}$ de l'homomorphisme $u:E\to A$ de $A$-modules. Cela prouve donc (i).

Pour établir (ii), notons que l'on peut supposer que $E$ est un $A$-module libre de rang $m$; puisque $\alpha_1^{\#}:E\otimes_A S_{(t)}\to(S(1))_{(t)}$ est surjectif et que $(S(1))_{(t)}$ est un $S_{(t)}$-module libre de rang $1$, la suite exacte
\[
  0\longrightarrow H_{(t)}\longrightarrow E\otimes_A S_{(t)}\xrightarrow{\alpha_1^{\#}}(S(1))_{(t)}\longrightarrow0
\]
est \emph{scindée}, et $H_{(t)}$ est donc un $S_{(t)}$-module projectif de rang $m-1$, d'où (ii).

Pour prouver (iii), examinons d'abord ce qui se passe au-dessus d'un point $x\in X-Y$. Alors $u:\mathscr E\to\sh{O}_X$, restreint à un voisinage de $x$, est surjectif, donc on peut se borner au cas où $\mathscr J=\sh{O}_X$, cas où $X'$ s'identifie canoniquement à $X$. Pour voir alors que le sous-préschéma fermé de $P$ image de $X'$ s'identifie à $Z$, on peut se borner comme au début au cas où $X=\Spec(A)$ est affine, et en outre $E=A^m$, de sorte que $S=A[T_0,\ldots,T_{m-1}]$. On peut en outre, avec les notations du début, supposer $u(t)\ne0$ (les valeurs de $t\in E$

